% Sección 7. Resume cifras ya citadas en las secciones P, 4 y 5.
\section{Conclusiones}
\label{sec:conclusiones}

\begin{enumerate}
  \item El pronóstico que pide el enunciado, a 30, 60 y 120 minutos, se equivoca a 60 minutos en \CPRutaModLX~\% de
        las subidas por ruta, \CPSentModLX~\% por sentido y \CPSegModLX~\% por segmento, en seis meses que no se
        usaron para elegir nada; el promedio por día y hora, en \CPRutaPromLX, \CPSentPromLX{} y \CPSegPromLX~\%. Lo que
        más aporta es lo observado en el mismo día: con solo lo que se sabe una semana antes, el mismo modelo falla
        \CPSemanalLX~\% por ruta. En los domingos corrientes no mejora al promedio.
  \item La demanda se pronostica con una semana de antelación con un WAPE de 13,6--13,9~\%, frente al 20,0~\% de la regla
        manual; la mejora sale de no aprender de datos corruptos y de combinar el LightGBM con un modelo fundacional.
  \item Lo que decide cuántos buses despachar es la elasticidad: con $\beta \approx 0{,}70$, recortar no paga con la
        contabilidad conservadora (el plan que recorta pierde 1.029 M COP al año; con la simétrica deja $+1.939$ M, a costa
        de pasajeros y sobrecarga). Con costos de fuente pública, el optimizador reasigna, y si eso paga depende
        de si los buses añadidos atraen pasajeros, algo que solo un piloto puede medir.
  \item El plan semanal queda mucho más cerca del óptimo que la operación real (1,05 frente a 2,80 salidas por hora de
        error); reoptimizar en el día no mejora.
  \item Visto persona a persona, el plan llena demasiado algunos buses: falta añadir el cupo por bus al optimizador antes
        de llevar el plan a la calle.
\end{enumerate}
